\section{Adding inheritance} \label{sec:inheritance}

 
The latest version of ShEx added the notion of shape \emph{inheritance} and was presented in \cite{ShExInheritance}:
a shape can \lean{extends} one or more other (abstract) labels,
and validation against an abstract label succeeds 
if the node matches \emph{some} concrete
descendant of it. 
Rather than retrofit this onto the 2017 formalization, 
we ported it as a genuine second formalism, 
\lean{ShexLean.Inheritance}, 
with its own \lean{ShapeExpr}/\lean{Schema}, 
reusing some common parts like  
(\lean{Graph}, \lean{Typing} and \lean{Label}).
The syntax itself grows in two ways beyond what 2017 needed: node
constraints become an arbitrary boolean predicate rather than a finite
\lean{ValueDescr}, and 
\lean{closed}/\lean{extra} shape modifiers and
\lean{extends} become first-class:

\begin{leanlisting}
structure Shape where
  closed : Bool
  extra  : Option (Finset IRI)
  triples : TripleExpr

inductive ShapeExpr where
  | ofShape (h : Shape) : ShapeExpr
  | extends (xs : Finset Label) (h : Shape) : ShapeExpr
  | constraint (c : NodeConstraint) : ShapeExpr
  | or (s1 s2 : ShapeExpr) : ShapeExpr
  | and (s1 s2 : ShapeExpr) : ShapeExpr
  | not (s : ShapeExpr) : ShapeExpr
\end{leanlisting}

A \lean{Shape}'s \lean{triples} field is itself a triple expression,
\lean{TripleExpr}, the inheritance formalism's counterpart to
\lean{TExpr} from Section~\ref{sec:shex_lean}: the empty expression, a
single triple constraint (property \lean{p}, target label \lean{z}), a
choice or sequence of two triple expressions, or a repeated triple
expression:

\begin{leanlisting}
inductive TripleExpr where
  | empty  : TripleExpr
  | tc (p : IRI) (z : Label) : TripleExpr
  | choice (e1 e2 : TripleExpr) : TripleExpr
  | seq (e1 e2 : TripleExpr) : TripleExpr
  | star (e : TripleExpr) : TripleExpr
\end{leanlisting}

The \lean{extends} case is what breaks the 2017 proof pattern at its foundation.
In the 2017 semantics, a shape reference is resolved purely by consulting
the typing, so \lean{Satisfies}/\lean{Matches} are
well-founded by ordinary structural recursion on the finite syntax tree.
An \lean{extends} clause instead recurses into \emph{other labels' own
definitions}, which implies a schema lookup,
not a subterm of the current expression,
so a plain structurally-recursive Lean definition
is not just inconvenient here,
it would be outright \emph{false} if the schema's extension hierarchy
happened to be cyclic.
The paper handles this with its own Lemma 1
(termination, proved separately from the semantics), so the Lean
development follows the same shape: \lean{Satisfies}
is first defined as a
total, \lean{Fuel}-indexed under-approximation. This is a standard technique
for giving a structurally-terminating definition to a recursion whose
termination is not otherwise obvious to Lean: an extra \lean{Fuel : ℕ}
argument is decremented on every recursive call, the recursion returns
\lean{False} outright once \lean{Fuel} reaches \lean{0}, and it is this
decrease in \lean{Fuel} — not the shape of the schema — that Lean accepts
as structural:

\begin{leanlisting}
def SatisfiesFuel : (Fuel : ℕ) → Schema → Graph → Typing → RDFNode →
    ShapeExpr → Prop
  | 0, _, _, _, _, _ => False
  ...
\end{leanlisting}

and a separate file, \lean{Termination.lean}, proves that for a
well-defined schema some large-enough,
canonical fuel value makes the answer stop depending on
\lean{Fuel} at all, recovering the paper's real,
fuel-free \lean{Satisfies}.
Turning the paper's one-sentence termination
argument ("the tree of recursive calls is finite due to the acyclic nature
of the extension hierarchy") into a Lean proof needs an actual rank
function on labels witnessing that acyclicity, combined into a single
measure that decreases on both kinds of recursive step:
ordinary AST descent and \lean{extends}-hops.

Definition 4 (correct typing) also gains a genuinely new case beyond a
direct port of the 2017 statement: an abstract label is satisfied not by
its own definition, but by the \emph{existence} of some concrete
descendant that is:

\begin{leanlisting}
def Schema.CorrectTyping (S : Schema) (G : Graph) (τ : Typing) : Prop :=
  ∀ n z, τ.has n z →
    (z ∉ S.A → Satisfies S G τ n (S.shapeDef z)) ∧
    (z ∈ S.A → ∃ x ∈ S.desc z \ S.A, Satisfies S G τ n (S.shapeDef x))
\end{leanlisting}

From there, Proposition 1, Lemma 2, and Algorithm 1 are structurally direct
ports of their 2017 counterparts — same per-stratum greatest-fixed-point
idea for the algorithm, same Fusion/Swap/Independence argument for Lemma 2
(whose own paper proof, again, is only sketched as "Proof. (Idea)...") —
but none of them transfer \emph{mechanically}. 
The 2017 monotonicity and schema-agreement lemmas these all lean on go through by plain structural
recursion, because 2017's \lean{Satisfies} never dereferences the schema.
Inheritance's \lean{Satisfies} does, via ancestor lookups, so every one of those
lemmas first had to be re-proved against the schema's dependency structure
and then generalized from a single stratification-tied restriction
(\lean{stratRestrict}) to an arbitrary dependency-closed label set
(\lean{Schema.restrictLabels}), before Lemma 2's schedule-agnostic argument and
Algorithm 1's per-stratum schema-swapping could reuse it. Algorithm 2
(\emph{prove}) has not yet been ported to this formalism.

In order to formalize the definitions we needed to tackle some potential
 places in which the paper didn't clearly specify all the conditions.  
For example, Table 3's \lean{extra} rule is
\emph{not monotone} in the typing without a further condition: 
growing the typing can move a triple out of the "extra-covered" bucket and 
into the main triple-expression match, which can break a bare 
(non-starred) triple constraint that needed an exact singleton. 
For instance, let node $n$ have two triples $t_1=(n,p,o_1)$ and
$t_2=(n,p,o_2)$, and let $h$ be a shape with \lean|closed := false|, 
\lean{extra := some {p}} and 
\lean{triples := p @z}.
Under $\tau_1=\{(o_1,z)\}$, $t_1$ matches \lean|p @z| and 
$t_2$ is covered by \lean|extra|, 
so $n$ satisfies $h$. Under $\tau_2=\tau_1\cup\{(o_2,z)\}$
both triples must match \lean|p @z|, which accepts exactly one, 
so $n$ no
longer satisfies $h$ although $\tau_1 \subseteq \tau_2$.
% We checked a concrete instance of this directly in Lean: 
% a shape with \lean{extra := some \{p\}} and
% \lean{triples := p @z}, satisfied under a typing $\tau_1$, provably
% \emph{stops} being satisfied under a strictly larger $\tau_2 \supseteq \tau_1$. 
The appendix proof of Proposition 1 in the paper states the fix for this: 
an invariant it calls "(*)": a shape's triple-constraint properties 
never overlap its own \lean{extra} set. 
That definition was not defined as a standalone definition 
in the main text: 
we named it explicitly as \lean{Schema.WellFormedExtra} and 
added it as a hypothesis to Proposition 1, Lemma 2 and Algorithm 1 of the inheritance formalism.

%\begin{table}[t]\centering\scriptsize
%\caption{Results formalized and gaps found.}\label{tab:summary}
%\begin{tabular}{lccl}
%\hline
%Result & 2017~\cite{BGP17} & Inheritance~\cite{ShExInheritance} & Gap found \\
%\hline
%Lemma 1 (neg.-free max.\ typing) & \checkmark & -- & -- \\
%Proposition 1 & \checkmark & \checkmark & closure (\lean{Closed}); (*) as \lean{WellFormedExtra} \\
%Lemma 2 & \checkmark & \checkmark & proof only sketched; full proof supplied \\
%Algorithm 1 (refine) & \checkmark & \checkmark & -- \\
%Algorithm 2 (prove) & \checkmark & not yet & -- \\
%Termination of $\vdash$ & n/a & \checkmark & rank function made explicit \\
%\hline
%\end{tabular}
%\end{table}

